\documentclass[14pt,usepdftitle=false,aspectratio=169]{beamer}
\usepackage{preambule}
\setbeamercolor{structure}{fg=black}
\usepackage[nopar]{bigoperators}\usepackage{polynomes,al}\newcommand{\corresp}[8]{\begin{array}[t]{@{}r@{}c@{}l@{}}#1&{}\longleftrightarrow{}&#2\\#4&{}\overset{\scriptscriptstyle#3}{\longmapsto}{}&#5\\#7&{}\overset{\scriptscriptstyle#6}{\longmapsfrom}{}&#8\end{array}}\DeclareMathOperator{\olddom}{dom}\newcommand{\dom}[1]{\olddom\l#1\r}\newcommand{\dashto}{\mathrel{-\mkern-6mu{\to}\mkern-22.5mu{\text{\color{white}\boldmath${\cdot}{\cdot}{\cdot}$}}\mkern1mu\phantom{.}}}
\hypersetup{pdftitle={Introduction à la géométrie algébrique -- Fonctions sur les variétés}}
\title{Introduction à la géométrie algébrique\\\emph{Fonctions sur les variétés}}
\author{}
\date{}
\begin{document}
\begin{frame}
    \titlepage
\end{frame}
\slideq{CNS topologiques pour qu'un ensemble algébrique $V$ soit irréductible}{1/14}
\slider{Pour tous $U_1$ et $U_2$ ouverts non vides de $V$, $U_1\cap U_2\neq\varnothing$\linebreak Tout ouvert de $V$ est dense}{1/14}
\slideq{Application rationnelle entre variétés affines $f\!:\!V\dashto W$}{2/14}
\slider{Fonction rationnelle $f\!:\!V\dashto\bbA_\bbk^n$ telle que $f\l\dom f\r\subset W$}{2/14}
\slideq{$\pol kV$\linebreak$V\subset\bbA_\bbk^n$ un ensemble algébrique}{3/14}
\slider{$\pol k{X_1,\cdots,X_n}/I\l V\r$}{3/14}
\slideq{Corps de fonctions sur une variété affine}{4/14}
\slider{$\fr kV=\operatorname{Frac}\l\pol kV\r$}{4/14}
\slideq{Anneau local de $V$ en $P$\linebreak$\calO_{V,P}$}{5/14}
\slider{$\set{f\in\pol kV,P\in\dom f}$\linebreak$\pol kV\lc\set{h^{-1},h\l P\r\neq0}\rc$}{5/14}
\slideq{$f\!:\!V\to W$ polynomiale est un isomorphisme}{6/14}
\slider{Il existe une fonction polynomiale $g\!:\!W\to V$ telle que $g\circ f=\id_V$ et $f\circ g=\id_W$}{6/14}
\slideq{Propriétés de $f^*$ pour $f\!:\!V\to W$}{7/14}
\slider{Si $f$ est une application polynomiale alors $\appl f{\pol kW}{\pol kV}g{g\circ f}$ est un morphisme de $\bbk$-algèbres\linebreak Tout morphisme de $\bbk$-algèbres $\varPhi\!:\!\pol kW\to\pol kV$ est de la forme $\varPhi=f^*$ pour $f\!:\!V\to W$ polynomiale\linebreak$\l g\circ f\r^*=f^*\circ g^*$}{7/14}
\slideq{$f\in\fr kV$ est régulière en $P\in V$}{8/14}
\slider{Il existe une représentation $f=\frac gh$ avec $h\l P\r\neq0$}{8/14}
\slideq{Propriétés de $\dom f$}{9/14}
\slider{$\dom f$ est ouvert dense pour la topologie de Zariski\linebreak Si $\bbk$ est algébriquement clos, alors $\dom f=V$ si et seulement si $f\in\pol kV$}{9/14}
\slideq{$f$ est une fonction polynomiale sur $V$}{10/14}
\slider{$f$ est la restriction à $V$ d'une application polynomiale $F\in\pol k{X_1,\cdots,X_n}$}{10/14}
\slideq{Variété affine sur un corps $\bbk$}{11/14}
\slider{Ensemble $V$ avec un anneau $\pol kV$ de fonction de $V$ dans $\bbk$ tels que $\pol kV$ est une $\bbk$-algèbre de type fini et, pour un choix de base $\l x_1,\cdots,x_n\r$ de $\pol kV$ sur $\bbk$, $\nappl{V}{\bbA_\bbk^n}{P}{\l x_1\l P\r,\cdots,x_n\l P\r\r}$ plonge $V$ dans un ensemble algébrique irréductible\linebreak Les variétés sont définies à isomorphisme près}{11/14}
\slideq{Fonction rationnelle $f\!:\!V\dashto\bbA_\bbk^n$\linebreak$V$ une variété affine}{12/14}
\slider{Application partielle définie par $f=\l f_1,\cdots,f_n\r$, $f_i\in\fr kV$\linebreak$\dom f=\bigcap{i=1}n{\dom{f_i}}$\linebreak $f$ est régulière en $P$ si $P\in\dom f$}{12/14}
\slideq{Correspondances induites par $I$ et $V$ pour les fonctions polynomiales}{13/14}
\slider{$\corresp{\set{J\leqslant\pol kV}}{\set{X\subset V}}V{I}{\set{\substack{\textstyle P\in V,\\\textstyle\forall f\in I,\\\textstyle f\l P\r=0}}}I{\set{\substack{\textstyle f\in\pol kV,\\\textstyle\forall P\in X,\\\textstyle f\l P\r=0}}}{X}$}{13/14}
\slideq{Domaine de définition de $f\in\fr kV$\linebreak$\dom f$}{14/14}
\slider{$\set{P\in V,\exists f=\frac gh,h\l P\r\neq0}$}{14/14}
\end{document}